\section{Lecture I: Springer fibers}

Springer fibers are classical and fundamental objects in geometric representation theory. Springer \cite{Springer} first discovered that their cohomology groups realized representations of Weyl groups, a phenomenon known as the Springer correspondence. As singular algebraic varieties, Springer fibers are interesting geometric objects by themselves. They are also connected to the representation theory of finite groups of Lie type via character sheaves.

\subsection{The setup}\label{ss:notation}  In this section, let $k$ be an algebraically closed field. Let $G$ be a connected reductive group over $k$ whose adjoint group is simple (so the adjoint group $G^{\ad}$ is determined by one of the seven series of Dynkin diagrams). Assume that $\chark$ is large compared to $G$.  Let $r$ be the rank of $G$.

Sometimes it will be convenient to fix a Cartan subalgebra $\frt$ of $\frg$, or equivalently a maximal torus $T\subset G$. Once we have done this, we may talk about the roots of the $T$-action on $\frg$.
 
Let $\sB$ be the flag variety of $G$: this is the $G$-homogeneous projective variety parametrizing Borel subgroups of $G$. Choosing a Borel subgroup $B\subset G$, we may identify $\sB$ with $G/B$.

Let $\frg$ be the Lie algebra of $G$.  For $X\in\frg$, let $C_{G}(X)$ denote the centralizer of $X$ in $G$, i.e., the stabilizer at $X$ of the adjoint action of $G$ on $\frg$. 

Let $\sN\subset\frg$ be the subvariety of nilpotent elements. This is a cone: it is stable under the action of $\Gm$ on $\frg$ by scaling.  It is known that there are finitely many $G$-orbits on $\sN$ under the adjoint action. 


\subsection{Springer fibers}

\subsubsection{The Springer resolution} The cotangent bundle $\tN:=T^{*}\sB$ classifies pairs $(e, B)$ where $e\in\sN$ and $B$ is a Borel subgroup of $G$ such that $e\in \frn$, where $\frn$ is the nilpotent radical of $\Lie B$. The {\em Springer resolution} is the forgetful map
\begin{equation*}
\pi:\tN\to\sN
\end{equation*}
sending $(e, B)$ to $e$. This map is projective.

For $e\in\sN$, the fiber $\sB_{e}:=\pi^{-1}(e)$ is called the {\em Springer fiber} of $e$. By definition, $\sB_{e}$ is the closed subscheme of $\sB$ consisting of those Borel subgroups $B\subset G$ such that $e$ is contained in the nilpotent radical of $\Lie B$. 

\subsubsection{The Grothendieck alteration}\label{sss:Gr alt}
Consider the variety $\tg$ of pairs $(X,B)$ where $X\in\frg$ and $B\in\sB$ such that $X\in\Lie B$. The forgetful map $(X,B)\mapsto X$
\begin{equation*}
\pi_{\frg}: \tg\to \frg
\end{equation*}  
is called the {\em Grothendieck alteration} \footnote{The term ``alteration'' refers to a proper, generically finite map whose source is smooth over $k$.}, also known as the {\em Grothendieck simultaneous resolution}. 
 
Let $\frc=\frg\sslash G\cong \frt\sslash W$ be the categorical quotient of $\frg$ by the adjoint action of $G$. Then we have a commutative diagram
\begin{equation}\label{tgg}
\xymatrix{\tg\ar[d]^{\pi_{\frg}}\ar[r]^{\wt\chi} & \frt\ar[d]\\
\frg \ar[r]^{\chi} & \frc}
\end{equation}
Here $\chi:\frg\to \frc$ is the natural quotient map and $\wt\chi:\tg\to \frt$ sends $(X,B)\in\tg$ to the image of $X$ in $\frb/\frn$ (where $\frb=\Lie B$ with nilpotent radical $\frn$), which can be canonically identified with $\frt$ (upon choosing a Borel containing $\frt$). The diagram \eqref{tgg} is Cartesian when restricting the left column to regular elements  \footnote{An element $X\in\frg$ is called {\em regular} if its centralizer in $G$ has dimension $r$, the rank of $G$.} in $\frg$. In particular, if we restrict the diagram \eqref{tgg} to the regular semisimple locus $\frc^{\rs}\subset \frc$   \footnote{An element $X\in\frt$ is {\em regular semisimple} if $\alpha(X)\neq0$ for any root $\alpha$. Denote by $\frt^{\rs}\subset\frt$ the open subscheme of regular semisimple elements in $\frt$. Since $\frt^{\rs}$ is stable under $W$, it is the full preimage of an open subset $\frc^{\rs}\subset\frc$. The open subset $\frg^{\rs}=\chi^{-1}(\frc^{\rs})$ is by definition the locus of {\em regular semisimple} elements of $\frg$.}, we see that $\pi_{\frg}|_{\tg^{\rs}}:\tg^{\rs}\to \frg^{\rs}$ is a $W$-torsor; i.e., there is an action of $W$ on $\pi_{\frg}^{-1}(\frg^{\rs})$ preserving the projection to $\frg^{\rs}$ making the fibers of $\pi_{\frg}$ into principal homogeneous spaces for $W$. The map $\pi_{\frg}$ becomes branched but still finite over the regular locus $\frg^{\reg}\subset \frg$.


For $X\in \frg$, we denote the fiber $\pi_{\frg}^{-1}(X)$ simply by $\tB_{X}$.  Restricting $\pi_{\frg}$ to $\sN$, the fibers $\tB_{e}$ for a nilpotent element $e\in\sN$ is the closed subvariety of $\sB$ consisting of those $B$ such that $e\in\Lie B$.  Clearly $\sB_{e}$ is a subscheme of $\tB_{e}$, and the two schemes have the same reduced structure. However, as schemes, $\tB_{e}$ and $\sB_{e}$ are different in general. See \S\ref{sss:regular Spr} below and Exercise \ref{exI:nonred Spr}.


